\chapter{Communication Complexity --- Pinsker}
%%% generated by the blueprint pipeline from TCSlib.CommunicationComplexity.NewmanTheorem.Pinsker
\section{Overview}

This module develops a self-contained proof of Pinsker's inequality, which bounds the
total variation distance between two probability measures by their KL divergence.
The argument proceeds via the Radon-Nikodym density and a variational characterisation
of total variation, culminating in
$2 \cdot \mathrm{TV}(\mu,\nu)^2 \le \mathrm{KL}(\mu \| \nu)$.

%%%%%%%%%%%%%%%%%%%%%%%%%%%%%%%%%%%%%%%%%%%%%%%%%%%%%%%%%%%%%%%%%%%%%%%%%%%%%%%
\section{Declarations}

\begin{definition}[Radon-Nikodym density]
\lean{CommunicationComplexity.rnDensity}
\label{CommunicationComplexity.rnDensity}
\leanok
For two probability measures $\mu$ and $\nu$ on a measurable space $\Omega$, the
\emph{Radon-Nikodym density} $\mathrm{rnDensity}(\mu,\nu)(x) \in \mathbb{R}$ is
the real part of the Radon-Nikodym derivative $\frac{d\mu}{d\nu}(x)$, obtained by
converting the extended non-negative real value $\mu.\mathrm{rnDeriv}\,\nu\,x$
to a real number via \texttt{ENNReal.toReal}.
\end{definition}

\begin{theorem}[Nonnegativity of the Radon–Nikodym density]
\lean{CommunicationComplexity.rnDensity_nonneg}
\label{CommunicationComplexity.rnDensity_nonneg}
\leanok
\uses{CommunicationComplexity.rnDensity}
\difficulty{1}
Let $\mu$ and $\nu$ be probability measures on a measurable space $\Omega$. Then at
every point $x \in \Omega$, the real-valued Radon–Nikodym density $\frac{d\mu}{d\nu}(x)$
satisfies $\frac{d\mu}{d\nu}(x) \ge 0$.
\end{theorem}

\begin{theorem}[Measurability of the Radon–Nikodym density]
\lean{CommunicationComplexity.measurable_rnDensity}
\label{CommunicationComplexity.measurable_rnDensity}
\leanok
\uses{CommunicationComplexity.rnDensity}
\difficulty{1}
Let $\mu$ and $\nu$ be probability measures on a measurable space $\Omega$, and consider
the real-valued Radon–Nikodym density $x \mapsto \left(\tfrac{d\mu}{d\nu}(x)\right)$,
obtained from the $[0,\infty]$-valued Radon–Nikodym derivative $\tfrac{d\mu}{d\nu}\colon
\Omega \to [0,\infty]$ by taking its real value at each point. Then this function
$\Omega \to \bbr$ is measurable.
\end{theorem}

\begin{theorem}[Integrability of the Radon–Nikodym density]
\lean{CommunicationComplexity.integrable_rnDensity}
\label{CommunicationComplexity.integrable_rnDensity}
\leanok
\uses{CommunicationComplexity.rnDensity}
\difficulty{1}
Let $\mu$ and $\nu$ be probability measures on a measurable space $\Omega$, and let
$f\colon \Omega \to \bbr$ be the real-valued Radon–Nikodym density of $\mu$ with respect
to $\nu$, that is, the function whose value at each point $x$ is the real number
obtained from the Radon–Nikodym derivative $\frac{d\mu}{d\nu}(x)$. Then $f$ is
integrable with respect to $\nu$.
\end{theorem}

\begin{theorem}[Integral of the Radon–Nikodym derivative equals one]
\lean{CommunicationComplexity.integral_rnDensity_eq_one_of_ac}
\label{CommunicationComplexity.integral_rnDensity_eq_one_of_ac}
\leanok
\uses{CommunicationComplexity.rnDensity}
\difficulty{2}
Let $\mu$ and $\nu$ be probability measures on a measurable space $\Omega$, and suppose
that $\mu$ is absolutely continuous with respect to $\nu$, written $\mu \ll \nu$. Then
the Radon–Nikodym derivative $\frac{d\mu}{d\nu}$ integrates to one against $\nu$:
\[
  \int_{\Omega} \frac{d\mu}{d\nu}(x)\,d\nu(x) = 1.
\]
\end{theorem}

\begin{theorem}[Integrability of the Radon–Nikodym density minus one]
\lean{CommunicationComplexity.integrable_rnDensity_sub_one}
\label{CommunicationComplexity.integrable_rnDensity_sub_one}
\leanok
\uses{CommunicationComplexity.integrable_rnDensity, CommunicationComplexity.rnDensity}
\difficulty{1}
Let $\mu$ and $\nu$ be probability measures on a measurable space $\Omega$, and let
$\frac{d\mu}{d\nu}\colon \Omega \to \bbr$ denote the real-valued Radon–Nikodym density
of $\mu$ with respect to $\nu$. Then the function $x \mapsto \frac{d\mu}{d\nu}(x) - 1$
is integrable with respect to $\nu$.
\end{theorem}

\begin{theorem}[Integrability of the deviation of the Radon–Nikodym density from one]
\lean{CommunicationComplexity.integrable_abs_rnDensity_sub_one}
\label{CommunicationComplexity.integrable_abs_rnDensity_sub_one}
\leanok
\uses{CommunicationComplexity.integrable_rnDensity_sub_one, CommunicationComplexity.rnDensity}
\difficulty{1}
Let $\mu$ and $\nu$ be probability measures on a measurable space $\Omega$, and let
$\frac{d\mu}{d\nu} : \Omega \to \bbr$ be the real-valued Radon–Nikodym density of $\mu$
with respect to $\nu$ (the extended nonnegative Radon–Nikodym derivative converted to a
real number). Then the function $x \mapsto \abs{\frac{d\mu}{d\nu}(x) - 1}$ is integrable
with respect to $\nu$.
\end{theorem}

\begin{theorem}[Mean of the Radon–Nikodym density minus one vanishes]
\lean{CommunicationComplexity.integral_rnDensity_sub_one_eq_zero_of_ac}
\label{CommunicationComplexity.integral_rnDensity_sub_one_eq_zero_of_ac}
\leanok
\uses{CommunicationComplexity.integrable_rnDensity, CommunicationComplexity.integral_rnDensity_eq_one_of_ac, CommunicationComplexity.rnDensity}
\difficulty{2}
Let $\mu$ and $\nu$ be probability measures on a measurable space $\Omega$, and suppose
that $\mu$ is absolutely continuous with respect to $\nu$. Writing $\frac{d\mu}{d\nu}$
for the Radon–Nikodym density of $\mu$ with respect to $\nu$ (the real-valued version of
the Radon–Nikodym derivative), one has \[ \int_{\Omega} \left(\frac{d\mu}{d\nu}(x) -
1\right)\, d\nu(x) = 0. \]
\end{theorem}

\begin{theorem}[Set difference of two measures as an integral of the density minus one]
\lean{CommunicationComplexity.measureReal_sub_eq_setIntegral_rnDensity_sub_one}
\label{CommunicationComplexity.measureReal_sub_eq_setIntegral_rnDensity_sub_one}
\leanok
\uses{CommunicationComplexity.integrable_rnDensity, CommunicationComplexity.rnDensity}
\difficulty{3}
Let $\mu$ and $\nu$ be probability measures on a measurable space $\Omega$ with $\mu$
absolutely continuous with respect to $\nu$ (that is, $\mu \ll \nu$), and let $S
\subseteq \Omega$. Writing $\frac{d\mu}{d\nu}$ for the (real-valued) Radon–Nikodym
derivative of $\mu$ with respect to $\nu$, one has
\[
  \mu(S) - \nu(S) \;=\; \int_{S} \left( \frac{d\mu}{d\nu}(x) - 1 \right) d\nu(x).
\]
\end{theorem}

\begin{theorem}[Positive part integrates to half the $L^1$ norm at mean zero]
\lean{CommunicationComplexity.integral_nonneg_part_eq_half_integral_abs_of_integral_eq_zero}
\label{CommunicationComplexity.integral_nonneg_part_eq_half_integral_abs_of_integral_eq_zero}
\leanok
\difficulty{4}
Let $\mu$ be a finite measure on a measurable space $\Omega$, and let $g : \Omega \to
\bbr$ be measurable and integrable with mean zero, $\int_\Omega g\,d\mu = 0$. Then the
integral of $g$ over the region where it is nonnegative equals half its $L^1$ norm:
\[
  \int_{\{x\,:\,g(x)\ge 0\}} g(x)\,d\mu \;=\; \tfrac{1}{2}\int_\Omega \abs{g(x)}\,d\mu.
\]
\end{theorem}

\begin{theorem}[Supremum of absolute set-integrals of a mean-zero function]
\lean{CommunicationComplexity.sSup_abs_setIntegral_eq_nonneg_part_of_integral_eq_zero}
\label{CommunicationComplexity.sSup_abs_setIntegral_eq_nonneg_part_of_integral_eq_zero}
\leanok
\difficulty{5}
Let $\mu$ be a finite measure on a measurable space $\Omega$, and let $g \colon \Omega
\to \bbr$ be measurable and integrable with mean zero, $\int_\Omega g\,d\mu = 0$. Then
the supremum of $\abs{\int_S g\,d\mu}$ taken over all measurable sets $S \subseteq
\Omega$ equals the integral of $g$ over the region where it is nonnegative:
\[
\sup_{S\text{ measurable}} \abs*{\int_S g\,d\mu} \;=\; \int_{\{x \,:\, g(x) \ge 0\}}
g\,d\mu.
\]
\end{theorem}

\begin{definition}[Density absolute integral]
\lean{CommunicationComplexity.densityAbsIntegral}
\label{CommunicationComplexity.densityAbsIntegral}
\leanok
\uses{CommunicationComplexity.rnDensity}
The quantity $\mathrm{densityAbsIntegral}(\mu,\nu) := \int |\mathrm{rnDensity}(\mu,\nu)(x) - 1|\,d\nu$ is the $L^1(\nu)$-distance of the Radon-Nikodym density from $1$.
\end{definition}

\begin{definition}[Density positive set]
\lean{CommunicationComplexity.densityPositiveSet}
\label{CommunicationComplexity.densityPositiveSet}
\leanok
\uses{CommunicationComplexity.rnDensity}
The \emph{density positive set} is $\{x \in \Omega : \mathrm{rnDensity}(\mu,\nu)(x) \ge 1\}$,
the region where $\mu$ locally dominates $\nu$.
\end{definition}

\begin{theorem}[Measurability of the density positive set]
\lean{CommunicationComplexity.measurableSet_densityPositiveSet}
\label{CommunicationComplexity.measurableSet_densityPositiveSet}
\leanok
\uses{CommunicationComplexity.densityPositiveSet, CommunicationComplexity.measurable_rnDensity}
\difficulty{1}
Let $\mu$ and $\nu$ be probability measures on a measurable space $\Omega$, and for each
$x \in \Omega$ let $\frac{d\mu}{d\nu}(x) \in \bbr$ denote the Radon–Nikodym derivative
of $\mu$ with respect to $\nu$, taken as a real number. Then the set $\left\{x \in
\Omega : \frac{d\mu}{d\nu}(x) \ge 1\right\}$, on which $\mu$ locally dominates $\nu$, is
a measurable subset of $\Omega$.
\end{theorem}

\begin{definition}[Density positive integral]
\lean{CommunicationComplexity.densityPositiveIntegral}
\label{CommunicationComplexity.densityPositiveIntegral}
\leanok
\uses{CommunicationComplexity.densityPositiveSet, CommunicationComplexity.rnDensity}
The \emph{density positive integral} is
$\mathrm{densityPositiveIntegral}(\mu,\nu) := \int_{\mathrm{densityPositiveSet}} (\mathrm{rnDensity}(\mu,\nu)(x)-1)\,d\nu$,
the excess mass of $\mu$ over $\nu$ on the region where $\mu$ dominates.
\end{definition}

\begin{theorem}[Total variation distance as the excess-mass integral]
\lean{CommunicationComplexity.tvDistanceSup_eq_densityPositiveIntegral_of_ac}
\label{CommunicationComplexity.tvDistanceSup_eq_densityPositiveIntegral_of_ac}
\leanok
\uses{CommunicationComplexity.densityPositiveIntegral, CommunicationComplexity.densityPositiveSet, CommunicationComplexity.integrable_rnDensity_sub_one, CommunicationComplexity.integral_rnDensity_sub_one_eq_zero_of_ac, CommunicationComplexity.measurable_rnDensity, CommunicationComplexity.measureReal_sub_eq_setIntegral_rnDensity_sub_one, CommunicationComplexity.rnDensity, CommunicationComplexity.sSup_abs_setIntegral_eq_nonneg_part_of_integral_eq_zero, CommunicationComplexity.tvDistanceSup}
\difficulty{4}
Let $\mu$ and $\nu$ be probability measures on a measurable space $\Omega$, and suppose
$\mu$ is absolutely continuous with respect to $\nu$. Write $\frac{d\mu}{d\nu}$ for the
Radon–Nikodym derivative of $\mu$ with respect to $\nu$. Then the total variation
distance, taken as the supremum
\[
\mathrm{TV}_{\sup}(\mu,\nu) = \sup\bigl\{\,\abs{\mu(S) - \nu(S)} : S \subseteq \Omega
\text{ measurable}\,\bigr\},
\]
equals the excess mass of $\mu$ over $\nu$ on the region where $\mu$ locally dominates:
\[
\mathrm{TV}_{\sup}(\mu,\nu) = \int_{\{x\,:\,\frac{d\mu}{d\nu}(x)\,\ge\,1\}}
\left(\frac{d\mu}{d\nu}(x) - 1\right)\,d\nu(x).
\]
\end{theorem}

\begin{theorem}[Total variation distance as an integral of density excess]
\lean{CommunicationComplexity.tvDistance_eq_densityPositiveIntegral_of_ac}
\label{CommunicationComplexity.tvDistance_eq_densityPositiveIntegral_of_ac}
\leanok
\uses{CommunicationComplexity.TVDistance.tvDistance_eq_tvDistanceSup, CommunicationComplexity.densityPositiveIntegral, CommunicationComplexity.tvDistance, CommunicationComplexity.tvDistanceSup_eq_densityPositiveIntegral_of_ac}
\difficulty{1}
Let $\mu$ and $\nu$ be probability measures on a measurable space $\Omega$, and suppose
that $\mu$ is absolutely continuous with respect to $\nu$. Write $\frac{d\mu}{d\nu}$ for
the (real-valued) Radon–Nikodym derivative of $\mu$ with respect to $\nu$. Then the
total variation distance between $\mu$ and $\nu$ equals the excess mass of $\mu$ over
$\nu$ on the region where $\mu$ locally dominates $\nu$:
\[
\tfrac{1}{2}\,\norm{\mu - \nu}_{\mathrm{TV}}
= \int_{\left\{\,x \in \Omega \;:\; \frac{d\mu}{d\nu}(x) \ge 1\,\right\}}
\left(\frac{d\mu}{d\nu}(x) - 1\right)\, d\nu.
\]
\end{theorem}

\begin{theorem}[Positive part equals half the total variation of the density]
\lean{CommunicationComplexity.densityPositiveIntegral_eq_half_densityAbsIntegral_of_ac}
\label{CommunicationComplexity.densityPositiveIntegral_eq_half_densityAbsIntegral_of_ac}
\leanok
\uses{CommunicationComplexity.densityAbsIntegral, CommunicationComplexity.densityPositiveIntegral, CommunicationComplexity.densityPositiveSet, CommunicationComplexity.integrable_rnDensity_sub_one, CommunicationComplexity.integral_nonneg_part_eq_half_integral_abs_of_integral_eq_zero, CommunicationComplexity.integral_rnDensity_sub_one_eq_zero_of_ac, CommunicationComplexity.measurable_rnDensity, CommunicationComplexity.rnDensity}
\difficulty{2}
Let $\mu$ and $\nu$ be probability measures on a measurable space $\Omega$, with $\mu$
absolutely continuous with respect to $\nu$, and write $\frac{d\mu}{d\nu}$ for the
Radon–Nikodym density of $\mu$ with respect to $\nu$. Then the excess mass of the
density above $1$, integrated over the region where the density is at least $1$, equals
half the $L^1(\nu)$-distance of the density from $1$:
\[
\int_{\{x\,:\,\frac{d\mu}{d\nu}(x)\,\ge\,1\}}
\left(\frac{d\mu}{d\nu}(x)-1\right)\,d\nu(x) \;=\; \frac{1}{2}\int_{\Omega}
\left|\frac{d\mu}{d\nu}(x)-1\right|\,d\nu(x).
\]
\end{theorem}

\begin{theorem}[A Fenchel-type inequality for the Kullback–Leibler function]
\lean{CommunicationComplexity.mul_le_klFun_add_exp_sub_one}
\label{CommunicationComplexity.mul_le_klFun_add_exp_sub_one}
\leanok
\difficulty{4}
For every real number $u \ge 0$ and every real number $y$,
\[
u\,y \;\le\; \left(u\ln u - u + 1\right) + e^{y} - 1,
\]
where the term $u\ln u$ is taken to be $0$ at $u = 0$.
\end{theorem}

\begin{theorem}[Exponential tilting by the cumulant generating function normalises to one]
\lean{CommunicationComplexity.integral_exp_sub_cgf_eq_one}
\label{CommunicationComplexity.integral_exp_sub_cgf_eq_one}
\leanok
\difficulty{3}
Let $\mu$ be a probability measure on a measurable space $\Omega$, let $X : \Omega \to
\bbr$ be a random variable, and let $t \in \bbr$. Write $\Lambda(t) = \log \int_\Omega
e^{t X(x)}\,d\mu(x)$ for the cumulant generating function of $X$ at $t$. If the function
$x \mapsto e^{t X(x)}$ is $\mu$-integrable, then
\[
\int_\Omega e^{\,t X(x) - \Lambda(t)}\,d\mu(x) = 1.
\]
\end{theorem}

\begin{theorem}[Integrability of the centred exponential moment]
\lean{CommunicationComplexity.integrable_exp_sub_cgf}
\label{CommunicationComplexity.integrable_exp_sub_cgf}
\leanok
\difficulty{2}
Let $(\Omega,\mu)$ be a measure space, let $X\colon\Omega\to\bbr$ be a measurable
function, and let $t\in\bbr$. Write $\Lambda(t)$ for the cumulant generating function of
$X$ under $\mu$, that is, the logarithm of the moment generating function evaluated at
$t$. If the function $x\mapsto e^{t\,X(x)}$ is $\mu$-integrable, then so is the function
$x\mapsto e^{t\,X(x)-\Lambda(t)}$.
\end{theorem}

\begin{theorem}[A variational inequality for the mean of a random variable]
\lean{CommunicationComplexity.variational_integral_mul_le_integral_klFun_add_cgf}
\label{CommunicationComplexity.variational_integral_mul_le_integral_klFun_add_cgf}
\leanok
\uses{CommunicationComplexity.integrable_exp_sub_cgf, CommunicationComplexity.integral_exp_sub_cgf_eq_one, CommunicationComplexity.mul_le_klFun_add_exp_sub_one}
\difficulty{5}
Let $\mu$ be a probability measure on a measurable space $\Omega$, let $f, X : \Omega
\to \bbr$, and let $t \in \bbr$. Suppose that $f \ge 0$ pointwise, that $\int_\Omega
f\,d\mu = 1$, and that $f$, the map $x \mapsto f(x)\log f(x) - f(x) + 1$, the map $x
\mapsto e^{tX(x)}$, and the product $fX$ are all $\mu$-integrable. Writing $\Lambda(t) =
\log \int_\Omega e^{tX}\,d\mu$ for the cumulant generating function of $X$ under $\mu$,
one has
\[
t \int_\Omega f(x)\,X(x)\,d\mu(x) \;\le\; \int_\Omega \bigl(f(x)\log f(x) - f(x) +
1\bigr)\,d\mu(x) + \Lambda(t).
\]
\end{theorem}

\begin{theorem}[A Pinsker-type inequality for the Kullback–Leibler generator]
\lean{CommunicationComplexity.half_integral_abs_sub_one_sq_le_integral_klFun}
\label{CommunicationComplexity.half_integral_abs_sub_one_sq_le_integral_klFun}
\leanok
\uses{CommunicationComplexity.variational_integral_mul_le_integral_klFun_add_cgf}
\difficulty{6}
Let $(\Omega, \mu)$ be a probability space, and let $\varphi(t) = t \log t - t + 1$ for
$t \ge 0$ denote the convex generator of the Kullback–Leibler divergence. Let $f :
\Omega \to \bbr$ be a measurable, integrable, nonnegative function with $\int_\Omega f
\, d\mu = 1$, so that $f$ is the density of a probability measure with respect to $\mu$,
and assume that $x \mapsto \varphi(f(x))$ is integrable with respect to $\mu$. Then
\[
\tfrac{1}{2}\Bigl(\int_\Omega \abs{f(x) - 1}\, d\mu\Bigr)^{2} \;\le\; \int_\Omega
\varphi(f(x))\, d\mu .
\]
\end{theorem}

\begin{theorem}[Total variation distance via the density's L¹-deviation from one]
\lean{CommunicationComplexity.tvDistance_eq_half_densityAbsIntegral_of_ac}
\label{CommunicationComplexity.tvDistance_eq_half_densityAbsIntegral_of_ac}
\leanok
\uses{CommunicationComplexity.densityAbsIntegral, CommunicationComplexity.densityPositiveIntegral_eq_half_densityAbsIntegral_of_ac, CommunicationComplexity.tvDistance, CommunicationComplexity.tvDistance_eq_densityPositiveIntegral_of_ac}
\difficulty{1}
Let $\mu$ and $\nu$ be probability measures on a measurable space $\Omega$, and suppose
that $\mu$ is absolutely continuous with respect to $\nu$. Then the total variation
distance between $\mu$ and $\nu$ equals one half of the density absolute integral, that
is,
\[
\mathrm{TV}(\mu,\nu) = \frac{1}{2}\int_{\Omega} \left| \frac{d\mu}{d\nu}(x) - 1 \right|
\, d\nu(x).
\]
\end{theorem}

\begin{theorem}[A Pinsker-type inequality for the Radon–Nikodym density]
\lean{CommunicationComplexity.half_densityAbsIntegral_sq_le_integral_klFun_rnDensity}
\label{CommunicationComplexity.half_densityAbsIntegral_sq_le_integral_klFun_rnDensity}
\leanok
\uses{CommunicationComplexity.densityAbsIntegral, CommunicationComplexity.half_integral_abs_sub_one_sq_le_integral_klFun, CommunicationComplexity.integrable_rnDensity, CommunicationComplexity.integral_rnDensity_eq_one_of_ac, CommunicationComplexity.measurable_rnDensity, CommunicationComplexity.rnDensity, CommunicationComplexity.rnDensity_nonneg}
\difficulty{2}
Let $\mu$ and $\nu$ be probability measures on a measurable space $\Omega$ with $\mu$
absolutely continuous with respect to $\nu$, and write $f = \frac{d\mu}{d\nu}$ for the
(real-valued) Radon–Nikodym density. Suppose the log-likelihood ratio $x \mapsto
\log\frac{d\mu}{d\nu}(x)$ is $\mu$-integrable. Then
\[
  \frac{1}{2}\left(\int_\Omega \abs{f(x) - 1}\,d\nu(x)\right)^{2}
  \;\le\;
  \int_\Omega \bigl(f(x)\log f(x) - f(x) + 1\bigr)\,d\nu(x),
\]
where the left-hand inner integral is the $L^1(\nu)$-distance of the density from $1$
and the right-hand integrand is the Kullback–Leibler integrand $t \mapsto t\log t - t +
1$.
\end{theorem}

\begin{theorem}[Two integral forms of the Kullback–Leibler divergence]
\lean{CommunicationComplexity.integral_klFun_rnDeriv_eq_kl_integral}
\label{CommunicationComplexity.integral_klFun_rnDeriv_eq_kl_integral}
\leanok
\uses{CommunicationComplexity.rnDensity}
\difficulty{1}
Let $\mu$ and $\nu$ be probability measures on a measurable space $\Omega$ with $\mu$
absolutely continuous with respect to $\nu$, and write $\frac{d\mu}{d\nu}$ for the
(real-valued) Radon–Nikodym density of $\mu$ with respect to $\nu$. Assume that the
log-likelihood ratio $x \mapsto \log\frac{d\mu}{d\nu}(x)$ is $\mu$-integrable. Then
\[
  \int_\Omega \varphi\!\left(\frac{d\mu}{d\nu}(x)\right) d\nu(x)
  \;=\;
  \int_\Omega \log\frac{d\mu}{d\nu}(x)\, d\mu(x),
\]
where $\varphi(t) = t\log t + 1 - t$.
\end{theorem}

\begin{theorem}[Pinsker-type bound via the KL integral]
\lean{CommunicationComplexity.density_l1_pinsker_le_kl_integral}
\label{CommunicationComplexity.density_l1_pinsker_le_kl_integral}
\leanok
\uses{CommunicationComplexity.densityAbsIntegral, CommunicationComplexity.half_densityAbsIntegral_sq_le_integral_klFun_rnDensity, CommunicationComplexity.integral_klFun_rnDeriv_eq_kl_integral, CommunicationComplexity.rnDensity}
\difficulty{2}
Let $\mu$ and $\nu$ be probability measures on a measurable space $\Omega$ with $\mu$
absolutely continuous with respect to $\nu$, and write $\frac{d\mu}{d\nu}$ for the
(real-valued) Radon–Nikodym density. Suppose the log-likelihood ratio $x \mapsto
\log\frac{d\mu}{d\nu}(x)$ is $\mu$-integrable. Then
\[
\frac{1}{2}\left(\int_{\Omega} \left|\frac{d\mu}{d\nu}(x) - 1\right|\,d\nu(x)\right)^{2}
  \;\le\;
  \int_{\Omega} \log\frac{d\mu}{d\nu}(x)\,d\mu(x),
\]
where the inner integral on the left is the $L^{1}(\nu)$-distance of the density from
the constant $1$.
\end{theorem}

\begin{theorem}[Pinsker's inequality under absolute continuity]
\lean{CommunicationComplexity.real_pinsker_inequality_of_ac_of_integrable}
\label{CommunicationComplexity.real_pinsker_inequality_of_ac_of_integrable}
\leanok
\uses{CommunicationComplexity.densityAbsIntegral, CommunicationComplexity.density_l1_pinsker_le_kl_integral, CommunicationComplexity.tvDistance, CommunicationComplexity.tvDistance_eq_half_densityAbsIntegral_of_ac}
\difficulty{2}
Let $\mu$ and $\nu$ be probability measures on a measurable space $\Omega$, and write
$\mathrm{TV}(\mu,\nu)$ for their total variation distance. Suppose $\mu$ is absolutely
continuous with respect to $\nu$, and that the log-likelihood ratio $x \mapsto
\log\frac{d\mu}{d\nu}(x)$ is $\mu$-integrable. Then
\[
  2\,\mathrm{TV}(\mu,\nu)^2 \;\le\; \int_\Omega \log\frac{d\mu}{d\nu}(x)\,d\mu(x),
\]
the right-hand side being the Kullback–Leibler divergence of $\mu$ from $\nu$.
\end{theorem}

\begin{theorem}[Pinsker's inequality under absolute continuity]
\lean{CommunicationComplexity.pinsker_inequality_of_ac_of_integrable}
\label{CommunicationComplexity.pinsker_inequality_of_ac_of_integrable}
\leanok
\uses{CommunicationComplexity.real_pinsker_inequality_of_ac_of_integrable, CommunicationComplexity.tvDistance}
\difficulty{2}
Let $\mu$ and $\nu$ be probability measures on a measurable space $\Omega$, and suppose
$\mu$ is absolutely continuous with respect to $\nu$ and that the log-likelihood ratio
$\log\frac{d\mu}{d\nu}$ is $\mu$-integrable. Then, as an inequality in the extended
nonnegative reals $\bbr_{\ge 0} \cup \{\infty\}$,
\[
  2\,\mathrm{TV}(\mu,\nu)^2 \;\le\; \mathrm{KL}(\mu \,\|\, \nu),
\]
where $\mathrm{TV}(\mu,\nu)$ is the total variation distance between $\mu$ and $\nu$ and
$\mathrm{KL}(\mu \,\|\, \nu)$ is the Kullback–Leibler divergence of $\mu$ from $\nu$.
\end{theorem}

\begin{theorem}[Pinsker's inequality under absolute continuity]
\lean{CommunicationComplexity.pinsker_inequality_of_ac}
\label{CommunicationComplexity.pinsker_inequality_of_ac}
\leanok
\uses{CommunicationComplexity.pinsker_inequality_of_ac_of_integrable, CommunicationComplexity.tvDistance}
\difficulty{2}
Let $\mu$ and $\nu$ be probability measures on a measurable space $\Omega$, and suppose
that $\mu$ is absolutely continuous with respect to $\nu$. Then the total variation
distance $\mathrm{TV}(\mu,\nu)$ and the Kullback–Leibler divergence $D_{\mathrm{KL}}(\mu
\,\|\, \nu)$ (both taking values in $[0,\infty]$) satisfy
\[
2\,\mathrm{TV}(\mu,\nu)^2 \le D_{\mathrm{KL}}(\mu \,\|\, \nu).
\]
\end{theorem}

\begin{theorem}[Pinsker's inequality]
\lean{CommunicationComplexity.pinsker_inequality}
\label{CommunicationComplexity.pinsker_inequality}
\leanok
\uses{CommunicationComplexity.pinsker_inequality_of_ac, CommunicationComplexity.tvDistance}
\difficulty{2}
Let $\mu$ and $\nu$ be probability measures on a measurable space $\Omega$, and let
$\mathrm{TV}(\mu,\nu)$ denote their total variation distance and
$\mathrm{KL}(\mu\,\|\,\nu)$ their Kullback–Leibler divergence. Then
\[
  2\,\mathrm{TV}(\mu,\nu)^2 \;\le\; \mathrm{KL}(\mu\,\|\,\nu),
\]
where the left-hand side is regarded as an element of the extended nonnegative reals
$[0,\infty]$.
\end{theorem}

\begin{theorem}[Pinsker's inequality for finite KL divergence]
\lean{CommunicationComplexity.two_mul_tvDistance_sq_le_toReal_klDiv}
\label{CommunicationComplexity.two_mul_tvDistance_sq_le_toReal_klDiv}
\leanok
\uses{CommunicationComplexity.pinsker_inequality, CommunicationComplexity.tvDistance}
\difficulty{2}
Let $\mu$ and $\nu$ be probability measures on a measurable space $\Omega$, and suppose
that the Kullback–Leibler divergence $\mathrm{KL}(\mu \,\|\, \nu)$ is finite. Then the
total variation distance $\mathrm{TV}(\mu,\nu) = \tfrac{1}{2}\|\mu -
\nu\|_{\mathrm{TV}}$ satisfies
\[
  2\,\mathrm{TV}(\mu,\nu)^2 \;\le\; \mathrm{KL}(\mu \,\|\, \nu).
\]
\end{theorem}
